\section{Articulación de la propuesta con las Misiones de la Política de Investigación e Innovación del Ministerio}


La presente propuesta se articula de manera directa con la misión “Soberanía Sanitaria y Bienestar Social” de la Política de Investigación e Innovación Orientada por Misiones (PIIOM), la cual busca fortalecer las capacidades del país para garantizar el acceso equitativo a servicios de salud, promover el bienestar de la población y fomentar el desarrollo de tecnologías estratégicas en el sector salud.

En este marco, el proyecto contribuye al desarrollo de un enfoque innovador para el apoyo al diagnóstico de trastornos mentales, mediante el análisis de señales EEG y el uso de modelos fundacionales con inferencia estocástica. Este enfoque responde a la necesidad de avanzar hacia sistemas de apoyo clínico más objetivos, reproducibles y confiables, reduciendo la dependencia exclusiva de evaluaciones subjetivas y mejorando la toma de decisiones en contextos de atención en salud mental.

Adicionalmente, la propuesta promueve el uso estratégico de datos biomédicos, incluyendo información no etiquetada, y al abordar problemáticas estructurales como la escasez de datos clínicos anotados, la variabilidad entre fuentes de información y la necesidad de modelos que generalicen adecuadamente en escenarios reales. En este sentido, el desarrollo de modelos robustos, generalizables y con estimaciones de incertidumbre contribuye a fortalecer la confiabilidad y seguridad de las herramientas de inteligencia artificial en salud, aspecto clave para su adopción en entornos clínicos.

Asimismo, el proyecto aporta al fortalecimiento de capacidades nacionales en inteligencia artificial aplicada a la salud, promoviendo la generación de conocimiento, el desarrollo tecnológico y la formación de talento humano altamente calificado. De igual forma, favorece la articulación entre actores del sistema de salud y el ecosistema de ciencia, tecnología e innovación, contribuyendo a la transferencia de conocimiento y a la apropiación social de tecnologías emergentes en el ámbito de la salud mental.